\scanpage{28}\mspage{22}
% printed form: blank music staves.
% @ top, continuation | text | (extension splittée). Sauf erreur
(extension splittée). Sauf erreur, ce groupe est le groupe de Weyl complété (par autom. ext.) de $D_4$ (le gpe des autom. ext. de $D_4$ correspondant à $\unsure{\mathfrak S|\mathscr S}_t$).

% @ upper third, dark interline | text | Donc cat. des syst.
Donc cat. des syst. de racines de type $D_4\simeq$ cat. des complexes cubiques munis d'un triangle

% @ upper third through centre | text | Précisons comment on peut associer
Précisons comment on peut associer un syst. de 24 racines (de type $D_4$) à un triangle $t$ dans $C$. Rappelons qu'à la situation $(C,t=\{u,v,w\})$ on associe \deleted{\unsure{des bij pour} $i,j\in t$, $i\ne j$, \unsure{entre}} \inserted{un syst. transitif d'isom} \deleted{\unsure{entre}} les ens $B(i)$ $(i\in t)$, où $B(i)\subset\unsure{\mathfrak P|\mathcal P}_2(\widetilde P_i)$ est formé des \deleted{couples} paires $(\tilde x,\tilde y)$ d'éléments non liés ; $\{x,y\}\in B(i)$ et $\{x',y'\}\in B(j)$ $(i\ne j)$ sont associés sss \deleted{\unsure{tt} $x\in$} $x,y$ \uline{sont liés} à $x',y'$
% @ middle, common set | text | Si on désigne par B
\deleted{\unsure{Soit B|Si $B$}} Si on désigne par $B$ l'ens \deleted{des} ``commun'' des $B(i)$, i.e. l'ens des parties \inserted{$b$} à 6 élts de $C$ mutuellement non liés (i.e. des bases) telles que $b\cap\underset{\substack{\simeq\\\widetilde P_i}}{\widehat E(i)^*}$ soit de cardinal 2 pour tt $i\in t$, on voit que \inserted{$\forall i\in t$} $B\overset{\sim}{\longrightarrow}B(i)$ par $b\longmapsto b\cap E(i)$, et $\operatorname{Card}B=\operatorname{Card}B(i)=\frac{8.6}{2}=24$.

% @ lower third, condition | text | La condition écrite sur b
\inserted{Je dis que} La condition écrite sur $b$ \deleted{\unsure{est équiv.}} équivaut à celle que $b\cap E(i)$ soit de cardinal 2 $\forall i\in t$, \deleted{car [il \deleted{faut vérifier} \inserted{suffit} que cette dernière implique que $b\cap t=\varnothing$ \unsure{ou on a} au plus \inserted{$\operatorname{card}$} $(b\cap t)=1$, puisque les pts de $t$ sont mut$^{\mathrm t}$ liés,} \inserted{car} si on avait $u\in t\cap b$, on aurait $\operatorname{Card} E(u)\cap b=\varnothing$] $b$ orthogonale à tout $i\in t$ (on dit que $b$ orthogonale à $t$).
% @ lower left, diagonal NB | text | Comme la somme
\margin{\uline{NB} Comme la somme $\sum_{i\in t}i=\unsure{\eta|\lambda}$ est orth. à \uline{tte} racine, les racines orth. à $t$ sont aussi celles orth. à \uline{deux} él. $i\ne j$ de $t$.}% slanted
% archive sheet number: 27
